Einstein's General Theory of Relativity predicts bending of light around
massive bodies. For simplicity, we assume that the bending of light happens at
a single point for each light ray, as shown in the figure. The angle of
bending, $\theta_{\mathrm{b}}$, is given by
\[ \theta_{\mathrm{b}} = \frac{2R_{\mathrm{sch}}}{r} \]
where $R_{\mathrm{sch}}$ is the Schwarzschild radius associated with that
gravitational body. We call $r$, the distance of the incoming light ray from
the parallel $x$-axis passing through the centre of the body, as the ``impact
parameter''.

\ioaafig{2016_TH_T10-f1.pdf}

A massive body thus behaves somewhat like a focusing lens. The light rays
coming from infinite distance beyond a massive body, and having the same
impact parameter $r$, converge at a point along the axis, at a distance $f_r$
from the centre of the massive body. An observer at that point will benefit
from huge amplification due to this gravitational focusing. The massive body
in this case is being used as a Gravitational Lensing Telescope for
amplification of distant signals.

\begin{description}
\item[(T10.1)] Consider the possibility of our Sun as a gravitational lensing
  telescope. Calculate the shortest distance, $f_{\mathrm{min}}$, from the
  centre of the Sun (in A.U.) at which the light rays can get focused.
  \hfill[6\,pt]
\item[(T10.2)] Consider a small circular detector of radius $a$, kept at a
  distance $f_{\mathrm{min}}$ centered on the $x$-axis and perpendicular to
  it. Note that only the light rays which pass within a certain annulus (ring)
  of width $h$ (where $h \ll R_\odot$) around the Sun would encounter the
  detector. The amplification factor at the detector is defined as the ratio
  of the intensity of the light incident on the detector in the presence of
  the Sun and the intensity in the absence of the Sun.

  Express the amplification factor, $A_{\mathrm{m}}$, at the detector in terms
  of $R_\odot$ and $a$. \hfill[8\,pt]
\item[(T10.3)] Consider a spherical mass distribution, such as dark matter in
  a galaxy cluster, through which light rays can pass while undergoing
  gravitational bending. Assume for simplicity that for the gravitational
  bending with impact parameter, $r$, only the mass $M(r)$ enclosed inside the
  radius $r$ is relevant.

  What should be the mass distribution, $M(r)$, such that the gravitational
  lens behaves like an ideal optical convex lens? \hfill[6\,pt]
\end{description}
